\documentclass[11pt,letterpaper]{amsart}
\usepackage[T1]{fontenc}
\usepackage{lmodern}
\usepackage{amsmath,amssymb,amsthm,mathtools}
\usepackage[margin=1.1in]{geometry}
\usepackage{enumitem}
\usepackage{needspace}
\usepackage{microtype}
\usepackage{xcolor}
\usepackage{hyperref}
\hypersetup{colorlinks=true,linkcolor=blue!45!black,urlcolor=blue!45!black,citecolor=blue!45!black,
 pdftitle={Invariant random C*-subalgebras of products of hyperbolic groups},
 pdfauthor={}}

\numberwithin{equation}{section}
\newtheorem{theorem}{Theorem}[section]
\newtheorem{proposition}[theorem]{Proposition}
\newtheorem{lemma}[theorem]{Lemma}
\newtheorem{corollary}[theorem]{Corollary}
\newtheorem*{theoremA}{Theorem~A}
\newtheorem*{corollaryB}{Corollary~B}
\theoremstyle{definition}
\newtheorem{definition}[theorem]{Definition}
\theoremstyle{remark}
\newtheorem{remark}[theorem]{Remark}

\newcommand{\Cr}{C_r^*}
\newcommand{\Sub}{\operatorname{Sub}}
\newcommand{\Prob}{\operatorname{Prob}}
\newcommand{\IRS}{\operatorname{IRS}}
\newcommand{\IRA}{\operatorname{IRA}}
\newcommand{\Ad}{\operatorname{Ad}}
\newcommand{\id}{\operatorname{id}}
\newcommand{\C}{\mathbb C}
\newcommand{\F}{\mathbb F}
\newcommand{\N}{\mathbb N}
\newcommand{\Z}{\mathbb Z}
\newcommand{\Ecyc}[2]{E_{\langle #2\rangle}^{(#1)}}
\newcommand{\Etrace}[1]{\mathbb E_{\tau_0^{#1}}}

\title[Random $C^*$-subalgebras of products]{Invariant random $C^*$-subalgebras\\
of products of hyperbolic groups}
\date{23 October 2026}
\keywords{Invariant random subalgebra, reduced group $C^*$-algebra,
hyperbolic group, boomerang, invariant random subgroup}

\begin{document}
\begin{abstract}
We present a proof of boomerang rigidity for finite products of
torsion-free nonamenable hyperbolic groups, extending the single-factor
theorem of Fima--Oppelmayer--Tao. Every boomerang unital subalgebra of
the reduced group $C^*$-algebra is a subgroup algebra. Consequently,
the subgroup-algebra correspondence induces an affine bijection
between invariant random subgroups and invariant random
$C^*$-subalgebras. The proof combines the selective averaging of
Amrutam--Jiang with recurrence in the space of subalgebras.
Coordinate expectations remove terms with a trivial active coordinate;
a common return sequence then isolates each nonzero Fourier
coefficient. The argument also applies to countable restricted direct
sums.
\end{abstract}
\maketitle

\section{Introduction}

Let $\Gamma$ be a countable discrete group. Conjugation by the left
regular representation defines an action on the space
$\Sub(\Cr(\Gamma))$ of unital $C^*$-subalgebras containing the ambient
unit. We equip this space with the Wijsman topology.
An invariant random $C^*$-subalgebra is an invariant Borel probability
measure on this space. Every invariant random subgroup gives such a
measure through
\[
 \iota_r:\Sub(\Gamma)\longrightarrow\Sub(\Cr(\Gamma)),
 \qquad H\longmapsto\Cr(H).
\]
Following Fima--Oppelmayer--Tao, a random subalgebra obtained in this
way is called \emph{geometric}.

Fima--Oppelmayer--Tao prove that every invariant random subalgebra is
geometric when $\Gamma$ is torsion-free, nonamenable, and hyperbolic
\cite[Corollary~6.9]{FOT}. Their stronger statement classifies
boomerang subalgebras, whose conjugates return to the original algebra
along a sequence of positive powers of each group element
\cite[Theorem~6.8]{FOT}. In the deterministic setting,
Amrutam--Jiang prove that every invariant unital subalgebra of
$\Cr(\Gamma_1\times\cdots\times\Gamma_n)$ is a normal subgroup algebra
when the factors have these hypotheses \cite[Theorem~5.4]{AJ}.

This paper extends the boomerang and random conclusions to finite
products of the same class of groups. The main statement is as follows.

\begin{theoremA}
Let $\Gamma=\Gamma_1\times\cdots\times\Gamma_n$, where $n\ge1$ and each
$\Gamma_i$ is a torsion-free nonamenable hyperbolic group.
If $B\in\Sub(\Cr(\Gamma))$ is boomerang, then
\[
 B=\Cr(H_B),\qquad
 H_B=\{g\in\Gamma:\lambda_g\in B\}.
\]
\end{theoremA}

\Needspace{9\baselineskip}
The probabilistic consequence includes a precise inverse to the
subgroup-algebra correspondence.

\begin{corollaryB}
For $\Gamma$ as in Theorem~A, every invariant random $C^*$-subalgebra
is geometric. The map
\[
 \IRS(\Gamma)\longrightarrow\IRA(\Cr(\Gamma)),
 \qquad \eta\longmapsto(\iota_r)_*\eta,
\]
is an affine bijection, with inverse
$\nu\mapsto(B\mapsto H_B)_*\nu$. It preserves and reflects ergodicity.
\end{corollaryB}

In particular, the conclusion applies to $\F_2\times\F_2$.
The subgroup $H_B$ need not split across the coordinates. For example, if
$\chi:\F_2\to\Z$ is onto, then
\[
 H=\{(g,h)\in\F_2\times\F_2:\chi(g)=\chi(h)\}
\]
is normal. If $\chi(s)=1$, then $\lambda_{(s,s)}\in\Cr(H)$, whereas
$\lambda_{(s,e)},\lambda_{(e,s)}\notin\Cr(H)$.
Thus the proof must recover a tuple unitary without requiring its
individual coordinate unitaries to belong to the algebra.

The proof follows the coefficient argument in
\cite[Lemmas~6.5--6.7 and Theorem 6.8]{FOT}.
The product argument requires two additional steps. First, simultaneous
selective averaging \cite[Theorem~3.3]{AJ} shows that coordinate cyclic
expectations preserve a boomerang algebra. Composing two such
expectations gives the coordinate trace map. Second, these trace maps
remove terms that are trivial in any nontrivial coordinate of a selected
group element. The reduced-word calculation then isolates the selected
coefficient, using one return sequence for a tuple of conjugating
elements.

\section{Notation and preliminaries}\label{sec:notation}

We use the notation of \cite{FOT} for subalgebras and subgroup
expectations. Put $A=\Cr(\Gamma)$, and let $\lambda_g$ denote the left
regular unitary associated with $g\in\Gamma$. The canonical trace is
$\tau_0$, and
\[
 \widehat a(g)=\tau_0(a\lambda_g^*)\qquad(a\in A,\ g\in\Gamma).
\]
For $B\in\Sub(A)$, write
\[
 d(a,B)=\inf_{b\in B}\|a-b\|.
\]
The Wijsman topology is the weakest topology making all these distance
functions continuous. The space $\Sub(A)$ is Polish when $A$ is
separable \cite[Theorem~2.3]{FOT}.

For $H\le\Gamma$, let $E_H:A\to\Cr(H)$ be the canonical conditional
expectation. Thus
\[
 E_H(\lambda_g)=1_H(g)\lambda_g.
\]
It is unital completely positive and contractive. For subgroups $H,K$,
\begin{equation}\label{eq:expectation-intersection}
 E_HE_K=E_KE_H=E_{H\cap K}.
\end{equation}
This follows on group unitaries and then by norm density.

Write $\alpha_g=\Ad(\lambda_g)$. Each $\alpha_g$ is a norm isometry and
induces a homeomorphism of $\Sub(A)$.
A subalgebra $B$ is \emph{boomerang} if, for every $s\in\Gamma$, there
are strictly increasing positive integers $m_j$ such that
$\alpha_{s^{m_j}}(B)\to B$ in the Wijsman topology.
We call $(m_j)$ a return sequence for $s$ at $B$.
For every $a\in B$ and such a sequence,
\begin{equation}\label{eq:return}
 d(\alpha_{s^{-m_j}}(a),B)
 =d(a,\alpha_{s^{m_j}}(B))\longrightarrow0.
\end{equation}

We recall two elementary distance estimates, also underlying
\cite[Lemma~6.7]{FOT}. If $T:A\to A$ is a linear contraction with
$T(B)\subseteq B$, then
\begin{equation}\label{eq:contraction}
 d(Ta,B)\le d(a,B).
\end{equation}
For $u,v\in A$,
\begin{equation}\label{eq:product-distance}
 d(uv,B)\le \|u\|d(v,B)+\|v\|d(u,B)+d(u,B)d(v,B).
\end{equation}
Indeed, choose $b,c\in B$ close to $u,v$ and use
$uv-bc=(u-b)v+b(v-c)$, together with
$\|b\|\le\|u\|+\|u-b\|$.

An element $t\ne e$ is called \emph{primitive} if $t=u^k$ implies
$k=\pm1$. In a torsion-free nonamenable hyperbolic group, every nonidentity
element has a primitive root, and its centralizer is infinite cyclic
with a primitive generator; see \cite[Remark~4.2]{AJ} and the proof
of \cite[Lemma~6.5]{FOT}. Every primitive element has a free partner:
there is $h$ for which $\langle t,h\rangle$ is free on $t,h$
\cite[Theorem~0.2]{AD}.

For products, all tensor products are minimal. Let $\tau_0^i$ be the
canonical trace on $\Cr(\Gamma_i)$. Following
\cite[Section~5.2]{AJ}, denote the coordinate trace expectation by
$\Etrace{i}$. On elementary tensors,
\[
 \Etrace{i}(a_1\otimes\cdots\otimes a_n)
 =a_1\otimes\cdots\otimes
       \bigl(\tau_0^i(a_i)1\bigr)\otimes\cdots\otimes a_n.
\]
For a primitive $t\in\Gamma_i$, we also use the coordinate notation
\[
 \Ecyc{i}{t}
 =\id\otimes\cdots\otimes E_{\langle t\rangle}
                       \otimes\cdots\otimes\id.
\]
The latter notation denotes the lift of the subgroup expectation in
factor $i$.

\section{Coordinate expectations}\label{sec:coordinates}

We use the following consequence of the selective averaging theorem
of Amrutam--Jiang.

\begin{lemma}[Selective averaging {\cite[Theorem~3.3 and Lemma 4.1]{AJ}}]
\label{lem:averaging}
Let $G$ be a torsion-free nonamenable hyperbolic group, let $t\in G$
be primitive, and let $(m_j)$ be a strictly increasing sequence of
positive integers. For every finite family $a_1,\ldots,a_N\in\Cr(G)$,
every $\varepsilon>0$, and every $J_0\ge1$, there is a finite convex
combination
\[
 \Phi=\sum_{\ell=1}^{M}\theta_\ell\alpha_{t^{m_{j_\ell}}},
 \qquad j_\ell\ge J_0,\quad
 \theta_\ell\ge0,\quad\sum_{\ell=1}^{M}\theta_\ell=1,
\]
such that
\[
 \|\Phi(a_k)-E_{\langle t\rangle}(a_k)\|<\varepsilon
 \qquad(1\le k\le N).
\]
The same assertion holds with $t$ replaced by $t^{-1}$.
\end{lemma}

The cited theorem gives equal-weight averages. The convex formulation
above suffices here. Appendix~\ref{app:averaging} recalls its proof.

Both $\Phi$ and $E_{\langle t\rangle}$ are unital completely positive:
the first is a convex combination of $*$-automorphisms, and the second
is a conditional expectation. Hence, for every unital $C^*$-algebra
$D$, their minimal tensor lifts $\Phi\otimes\id_D$ and
$E_{\langle t\rangle}\otimes\id_D$ are unital completely positive
and contractive. If
$a\in\Cr(G)\otimes_{\min}D$ and
$p=\sum_{k=1}^{N}a_k\otimes d_k$, then
\begin{align}
 &\|(\Phi\otimes\id_D)(a)
       -(E_{\langle t\rangle}\otimes\id_D)(a)\|\notag\\
 &\qquad\le
 2\|a-p\|+\sum_{k=1}^{N}
       \|\Phi(a_k)-E_{\langle t\rangle}(a_k)\|\,\|d_k\|.
 \label{eq:tensor}
\end{align}
Thus one simultaneous estimate on finitely many factor elements gives
the required norm approximation for each fixed tensor-product element.

\begin{proposition}\label{prop:cyclic}
Let $\Gamma=\Gamma_1\times\cdots\times\Gamma_n$ be as in Theorem~A.
If $B\in\Sub(\Cr(\Gamma))$ is boomerang, then
\[
 \Ecyc{i}{t}(B)\subseteq B
\]
for every $i$ and every primitive $t\in\Gamma_i$.
\end{proposition}
\begin{proof}
Fix $a\in B$ and embed $t$ in coordinate $i$, writing $\widetilde t$
for the resulting element of $\Gamma$. Choose a return sequence
$(m_j)$ for $\widetilde t$ at $B$. Equation~\eqref{eq:return} gives
\[
 d(\alpha_{\widetilde t^{-m_j}}(a),B)\longrightarrow0.
\]
Fix $\varepsilon>0$ and choose $J_0$ so that these distances are below
$\varepsilon$ for $j\ge J_0$.

Identify $A$ with $\Cr(\Gamma_i)\otimes_{\min}D$, where $D$ is the
tensor product of the other factors. Choose
$p=\sum_{k=1}^{N}a_k\otimes d_k$ with $\|a-p\|<\varepsilon$, and put
$M_0=\sum_k\|d_k\|$.
Apply Lemma~\ref{lem:averaging} to $t^{-1}$ and the tail $j\ge J_0$,
with error $\varepsilon/(1+M_0)$ on each $a_k$.
The resulting map $\Psi=\Phi\otimes\id_D$ is a convex combination
of the conjugations $\alpha_{\widetilde t^{-m_j}}$ from that tail.
By \eqref{eq:tensor},
\[
 \|\Psi(a)-\Ecyc{i}{t}(a)\|<3\varepsilon.
\]
Every summand of $\Psi(a)$ is within $\varepsilon$ of $B$. Since $B$
is convex, $d(\Psi(a),B)\le\varepsilon$. It follows that
$d(\Ecyc{i}{t}(a),B)<4\varepsilon$.
Letting $\varepsilon\to0$ and using closedness of $B$ proves the assertion.
\end{proof}

\begin{corollary}\label{cor:expectations}
Under the hypotheses of Proposition~\ref{prop:cyclic},
$\Etrace{i}(B)\subseteq B$ for every $i$.
Moreover,
\[
 E_{C_\Gamma(w)}(B)\subseteq B\qquad(w\in\Gamma).
\]
If $J\subseteq\{1,\ldots,n\}$ and $t_i\in\Gamma_i$ is primitive
for $i\in J$, then the expectation onto
\[
 \Lambda=
 \prod_{i\in J}\langle t_i\rangle
 \times\prod_{i\notin J}\{e\}
\]
also preserves $B$.
\end{corollary}
\begin{proof}
Choose a primitive $u\in\Gamma_i$ and a free partner $h$.
Then $v=huh^{-1}$ is primitive and
$\langle u\rangle\cap\langle v\rangle=\{e\}$.
Equation~\eqref{eq:expectation-intersection} gives
\begin{equation}\label{eq:coordinate-trace}
 \Etrace{i}=\Ecyc{i}{u}\Ecyc{i}{v}.
\end{equation}
Proposition~\ref{prop:cyclic} proves the first assertion.

For $w=(w_1,\ldots,w_n)$,
\[
 C_\Gamma(w)=\prod_{i=1}^{n}C_{\Gamma_i}(w_i).
\]
If $w_i\ne e$, the corresponding centralizer is generated by a
primitive element; if $w_i=e$, it is the whole factor.
Thus $E_{C_\Gamma(w)}$ is a composition of coordinate cyclic
expectations and identity maps.
Finally,
\begin{equation}\label{eq:lambda-expectation}
 E_\Lambda
 =\left(\prod_{i\in J}\Ecyc{i}{t_i}\right)
  \left(\prod_{i\notin J}\Etrace{i}\right).
\end{equation}
All these maps commute by \eqref{eq:expectation-intersection}.
\end{proof}

\section{Mixed Fourier coefficients}\label{sec:coefficients}

The coefficient argument begins with the reduced-word calculation in
\cite[Lemma~2.3]{AJ} and \cite[Lemma~6.6]{FOT}.

\begin{lemma}\label{lem:words}
Let $t,q$ freely generate a free subgroup of a group $G$.
For $r,k,l\in\Z\setminus\{0\}$, put
\[
 w=t^rqt^rq^{-1},\qquad z=t^kqt^lq^{-1}.
\]
Then $z\in C_G(w)$ if and only if $k=l=r$.
\end{lemma}
\begin{proof}
The products $wz$ and $zw$ are reduced alternating words in
$\langle t\rangle*\langle q\rangle$. Equality of their first
$t$-syllables gives $k=r$, and equality of their next $t$-syllables
gives $l=r$. The converse is immediate.
Both products lie in the embedded free subgroup, so this comparison
also tests commutation in $G$.
\end{proof}

Fix a nonempty $J\subseteq\{1,\ldots,n\}$, primitive $t_i\in\Gamma_i$,
and nonzero integers $r_i$ for $i\in J$. Set
\[
 \Lambda=\prod_{i\in J}\langle t_i\rangle
                \times\prod_{i\notin J}\{e\}.
\]
For $\boldsymbol k=(k_i)_{i\in J}\in\Z^J$, write $t^{\boldsymbol k}$
for the element with coordinates $t_i^{k_i}$ on $J$ and identity
elsewhere. Put $g=t^{\boldsymbol r}$.
Choose $q_i\in\Gamma_i$ free from $t_i$, and let $q$ have coordinates
$q_i$ on $J$ and identity elsewhere.

\begin{lemma}\label{lem:coefficient}
Let $x\in\Cr(\Lambda)$ and $c=\widehat x(g)$, and put $w=gqgq^{-1}$.
If $\Etrace{i}(x)=0$ for every $i\in J$, then
\begin{equation}\label{eq:isolation}
 E_{C_\Gamma(w)}(x\lambda_qx\lambda_q^*)=c^2\lambda_w.
\end{equation}
For every $x\in\Cr(\Lambda)$, without the preceding vanishing
assumption,
\begin{equation}\label{eq:recovery}
 E_\Lambda(\lambda_w\lambda_qx^*\lambda_q^*)
 =\overline c\,\lambda_g.
\end{equation}
\end{lemma}
\begin{proof}
First let
\[
 p=\sum_{\boldsymbol k\in(\Z\setminus\{0\})^J}
        p_{\boldsymbol k}\lambda_{t^{\boldsymbol k}}
\]
have finite support. In $p\lambda_qp\lambda_q^*$, the ordered pair
$(\boldsymbol k,\boldsymbol l)$ contributes
$p_{\boldsymbol k}p_{\boldsymbol l}$ at the tuple with coordinates
$t_i^{k_i}q_it_i^{l_i}q_i^{-1}$ on $J$.
Such a tuple belongs to $C_\Gamma(w)$ exactly when each coordinate
commutes with $t_i^{r_i}q_it_i^{r_i}q_i^{-1}$.
Lemma~\ref{lem:words} forces
$\boldsymbol k=\boldsymbol l=\boldsymbol r$.
Consequently
\[
 E_{C_\Gamma(w)}(p\lambda_qp\lambda_q^*)
 =p_{\boldsymbol r}^{\,2}\lambda_w.
\]

To pass from polynomials to $x$, define
\begin{equation}\label{eq:Q}
 Q_J=\prod_{i\in J}(\id-\Etrace{i}).
\end{equation}
The coordinate trace maps commute. The hypotheses give $Q_Jx=x$,
and $\|Q_J\|\le2^{|J|}$.
Take $p_m\in\C[\Lambda]$ converging to $x$ in norm and set
$z_m=Q_Jp_m$. Then $z_m\to x$ in norm, and each $z_m$ has only
nonzero exponents in every coordinate of $J$.
The estimate
\[
 \|z_m\lambda_qz_m\lambda_q^*-x\lambda_qx\lambda_q^*\|
 \le(\|z_m\|+\|x\|)\|z_m-x\|
\]
and continuity of Fourier coefficients prove \eqref{eq:isolation}.
This passage uses polynomial density, without requiring convergence
of Fourier partial sums.

For the second identity, freeness gives
$\Lambda\cap q\Lambda q^{-1}=\{e\}$.
Checking group unitaries and using norm density yields
\[
 E_\Lambda(\lambda_qy\lambda_q^*)=\tau_0(y)1
 \qquad(y\in\Cr(\Lambda)).
\]
Since $\lambda_w=\lambda_g\lambda_q\lambda_g\lambda_q^*$,
bimodularity of $E_\Lambda$ gives
\[
 \begin{aligned}
 E_\Lambda(\lambda_w\lambda_qx^*\lambda_q^*)
 &=\lambda_g
     E_\Lambda(\lambda_q(\lambda_gx^*)\lambda_q^*)\\
 &=\tau_0(\lambda_gx^*)\lambda_g
 =\overline c\,\lambda_g .
 \end{aligned}
\]
Only membership of the bimodule factor $\lambda_g$ in $\Cr(\Lambda)$
is used here; its membership in $B$ is the conclusion still to be proved.
\end{proof}

\begin{remark}
The coordinate vanishing assumptions in \eqref{eq:isolation} remove
the case $k_i=l_i=0$, which otherwise contributes an identity in
coordinate $i$. The coefficients of $p$ need not factor across
coordinates. The surviving coefficient is the square of the single
coefficient $p_{\boldsymbol r}$.
\end{remark}

\section{Proof of the main statements}\label{sec:proof}

\begin{proof}[Proof of Theorem~A]
Put $H_B=\{g\in\Gamma:\lambda_g\in B\}$.
This is a subgroup and $\Cr(H_B)\subseteq B$.
We prove that
\begin{equation}\label{eq:extraction}
 a\in B,\quad \widehat a(g)\ne0
 \quad\Longrightarrow\quad \lambda_g\in B.
\end{equation}

The case $g=e$ follows from unitality.
Otherwise let $J=\{i:g_i\ne e\}$.
For each $i\in J$, write $g_i=t_i^{r_i}$ with $t_i$ primitive
and $r_i\ne0$. Form $\Lambda$ and $Q_J$ as in
Section~\ref{sec:coefficients}, and set
\begin{equation}\label{eq:centered-element}
 x=Q_JE_\Lambda(a).
\end{equation}
Corollary~\ref{cor:expectations} implies $x\in B\cap\Cr(\Lambda)$.
Moreover $\Etrace{i}(x)=0$ for every $i\in J$, and
\[
 c:=\widehat x(g)=\widehat a(g)\ne0.
\]
Indeed, $E_\Lambda$ retains the coefficient at $g$, while
$\Etrace{i}$ annihilates that coefficient for each $i\in J$.

Choose $h_i$ free from $t_i$ for every $i\in J$, and let $h\in\Gamma$
have these coordinates and identity elsewhere.
Choose a single return sequence $(m_j)$ for $h$ at $B$.
Put
\[
 q_j=h^{-m_j},\qquad
 w_j=gq_jgq_j^{-1},\qquad
 \delta_j=d(\lambda_{q_j}x\lambda_{q_j}^*,B).
\]
Equation~\eqref{eq:return} gives $\delta_j\to0$.
For every $j$ and $i\in J$, the pair $t_i,h_i^{-m_j}$ is free:
substitution into a reduced word in these two elements leaves a
nontrivial reduced word in $t_i,h_i$.
Thus Lemma~\ref{lem:coefficient} applies with $q=q_j$.

Let $K=\|x\|$. Since $x\in B$,
\[
 d(x\lambda_{q_j}x\lambda_{q_j}^*,B)\le K\delta_j.
\]
The expectation $E_{C_\Gamma(w_j)}$ is contractive and preserves $B$
by Corollary~\ref{cor:expectations}. Equation~\eqref{eq:isolation}
therefore gives
\begin{equation}\label{eq:approximate-unitary}
 \varepsilon_j:=d(\lambda_{w_j},B)
 \le\frac{K}{|c|^2}\delta_j\longrightarrow0.
\end{equation}
Adjoints preserve distance to $B$. Using
\eqref{eq:product-distance}, \eqref{eq:contraction}, and
\eqref{eq:recovery}, we obtain
\[
 \begin{aligned}
 |c|\,d(\lambda_g,B)
 &=d\bigl(E_\Lambda(
       \lambda_{w_j}\lambda_{q_j}x^*\lambda_{q_j}^*),B\bigr)\\
 &\le\delta_j+K\varepsilon_j+\varepsilon_j\delta_j
 \longrightarrow0.
 \end{aligned}
\]
This proves \eqref{eq:extraction}.
The same return sequence controls all active coordinates throughout.

For $a\in B$, \eqref{eq:extraction} shows that its Fourier coefficients
vanish outside $H_B$. Hence $a-E_{H_B}(a)$ has all Fourier coefficients
zero. Such an element vanishes on $\delta_e$ and, by commutation with
the right regular representation, on every $\delta_g$.
It is therefore zero. Thus $a=E_{H_B}(a)\in\Cr(H_B)$, proving the
opposite inclusion.
\end{proof}

\begin{proof}[Proof of Corollary~B]
Let $\nu$ be an invariant random subalgebra of $\Cr(\Gamma)$.
By \cite[Proposition~6.3]{FOT}, almost every $B$ is boomerang.
This is Poincar\'e recurrence applied to a countable base of
$\Sub(\Cr(\Gamma))$, followed by a countable intersection over
$s\in\Gamma$. Theorem~A gives $\nu(I_\Gamma)=1$, where
\[
 I_\Gamma=\iota_r(\Sub(\Gamma)).
\]

For completeness, the subgroup recovery map is Borel, since
\[
 \{B:g\in H_B\}=\{B:d(\lambda_g,B)=0\}
\]
is Wijsman closed for each $g$.
It is equivariant because
$H_{\lambda_sB\lambda_s^*}=sH_Bs^{-1}$.
The map $\iota_r$ is a Borel isomorphism onto $I_\Gamma$
\cite[Theorem~2.13]{FOT}. Also $H_{\Cr(H)}=H$:
if $\lambda_g\in\Cr(H)$, applying $E_H$ shows that $g\in H$.

Set $\eta=(B\mapsto H_B)_*\nu$.
Equivariance makes $\eta$ an IRS, and Theorem~A yields
$(\iota_r)_*\eta=\nu$. Applying the recovery map to any other
representing measure proves uniqueness. Both pushforward maps are
affine. Their equivariant Borel inverse correspondence identifies
invariant measurable sets, and hence preserves and reflects
ergodicity.
\end{proof}

\section{Restricted direct sums and scope}\label{sec:extensions}

\begin{corollary}\label{cor:restricted-sum}
Let $\Gamma=\bigoplus_{i\in\N}\Gamma_i$, where each $\Gamma_i$ is a
torsion-free nonamenable hyperbolic group.
Every boomerang unital subalgebra of $\Cr(\Gamma)$ is a subgroup
algebra. Every invariant random subalgebra is geometric, with the
unique IRS given by the map $B\mapsto H_B$.
\end{corollary}
\begin{proof}
For each $i$,
\[
 \Cr(\Gamma)=\Cr(\Gamma_i)\otimes_{\min}
                   \Cr\left(\bigoplus_{k\ne i}\Gamma_k\right).
\]
The proof of Proposition~\ref{prop:cyclic} applies to this decomposition.
Equation~\eqref{eq:coordinate-trace} then proves preservation of all
coordinate trace expectations.

Fix a finite set $J\subseteq\N$ and put
$\Gamma_J=\prod_{i\in J}\Gamma_i$.
Choose finite sets $F_m$ increasing to $\N\setminus J$.
The maps $\prod_{i\in F_m}\Etrace{i}$ converge in norm on each element
to $E_{\Gamma_J}$. They eventually agree with $E_{\Gamma_J}$ on
each group polynomial, since every element of $\Gamma$ has finite
coordinate support. Contractivity gives convergence on all of
$\Cr(\Gamma)$. Closedness of $B$ implies
$E_{\Gamma_J}(B)\subseteq B$.

For a selected nonidentity $g$, its set $J$ of nontrivial coordinates
is finite. Apply $E_{\Gamma_J}$, the cyclic expectations on $J$, and
$Q_J$ to construct \eqref{eq:centered-element}.
The coefficient calculation and the return sequence use only these
coordinates. Each expectation $E_{C_\Gamma(w_j)}$ is a composition
of the finitely many cyclic expectations on $J$, with the identity
on the remaining coordinates.
The proof of \eqref{eq:extraction} is therefore unchanged.
Subgroup recovery and the proof for invariant measures apply because
$\Gamma$ is countable.
\end{proof}

The factor properties used in the argument can be stated independently
of hyperbolicity:
\begin{enumerate}[label=(\roman*),leftmargin=*]
\item every nonidentity element has an infinite cyclic centralizer
with a primitive generator;
\item every primitive element admits a free partner;
\item simultaneous selective averaging onto the primitive cyclic
subgroup is available along every subsequence of powers, in the form
of Lemma~\ref{lem:averaging}.
\end{enumerate}
The same proof applies to finite products and countable restricted
direct sums of nontrivial countable groups with these properties.
The tensor estimate \eqref{eq:tensor} requires no nuclearity or
approximation property.

These hypotheses have not been established here for general
torsion-free acylindrically hyperbolic groups. In particular, cyclic
centralizers are not available for all their nonidentity elements.
No assertion is made for unrestricted infinite direct products.

The classification reduces the study of invariant random subalgebras
for these groups to the study of their IRSs. It does not enumerate
those IRSs, nor does it resolve the general
questions relating $C^*$-ISR, von Neumann ISR, and $C^*$-simplicity in
\cite[Questions~6.11--6.12]{AJ}.

\appendix
\section{Selective averaging along subsequences}\label{app:averaging}

We record the convex averaging argument used in
Lemma~\ref{lem:averaging}. It is adapted from the boundary-state
argument of \cite[Lemma~3.1 and Proposition 3.2]{AJ}.

\Needspace{8\baselineskip}
\begin{lemma}\label{lem:boundary}
Let $G$ be a countable group acting continuously on a compact
metrizable space $X$. Suppose that $t\in G$ has north--south dynamics
in both time directions, with fixed points $t^+,t^-$, and
\[
 s\{t^+,t^-\}\cap\{t^+,t^-\}=\varnothing
 \qquad(s\notin\langle t\rangle).
\]
Then the conclusion of Lemma~\ref{lem:averaging} holds for $t$.
\end{lemma}
\begin{proof}
Put $F=\{t^+,t^-\}$.
Extend a state $\varphi$ on $\Cr(G)$ to a state $\omega$ on
$C(X)\rtimes_rG$, using the canonical unital inclusion.
If $\rho=\omega|_{C(X)}$, then
\[
 (\omega\circ\alpha_{t^{-n}})(f)
 =\int_X f(t^nx)\,d\rho(x)\qquad(f\in C(X)).
\]
North--south dynamics and dominated convergence show that these
restrictions converge to a probability measure supported on $F$.

Let $\omega_\infty$ be a weak-* cluster point of the states
$\omega\circ\alpha_{t^{-n}}$.
For $s\notin\langle t\rangle$, choose $f\in C(X)$ with
$0\le f\le1$, $f|_F=1$, and $f|_{sF}=0$.
Cauchy--Schwarz gives
\begin{align*}
 |\omega_\infty(f\lambda_s)|^2
 &\le\omega_\infty(f)
       \omega_\infty(\lambda_s^*f\lambda_s)=0,\\
 |\omega_\infty((1-f)\lambda_s)|^2
 &\le\omega_\infty(1-f)
       \omega_\infty(\lambda_s^*(1-f)\lambda_s)=0.
\end{align*}
Thus $\omega_\infty(\lambda_s)=0$.
Conjugation by powers of $t$ fixes $\Cr(\langle t\rangle)$ pointwise.
Consequently, the restriction of every such cluster state to $\Cr(G)$
is $\varphi\circ E_{\langle t\rangle}$: the functionals agree on
group unitaries, and hence on all of $\Cr(G)$ by norm density.
Weak-* compactness therefore gives
\[
 (\omega\circ\alpha_{t^{-n}})|_{\Cr(G)}
 =\varphi\circ\alpha_{t^{-n}}
 \xrightarrow{\;w^*\;}\varphi\circ E_{\langle t\rangle}.
\]
This argument establishes convergence of the restricted states;
convergence of the full states on the crossed product is not needed.

Every bounded linear functional is a linear combination of states.
Therefore $\alpha_{t^{-n}}(a)$ converges to
$E_{\langle t\rangle}(a)$ in the Banach weak topology for every
$a\in\Cr(G)$.
For a finite family, the corresponding tuple converges weakly in a
finite direct sum of copies of $\Cr(G)$.
By Hahn--Banach separation, its limit belongs to the norm-closed convex
hull of every subsequence tail. This gives one common convex
combination for the family.
Interchanging $t$ and $t^{-1}$ proves the assertion in both directions.
\end{proof}

For primitive elements of torsion-free nonamenable hyperbolic groups,
the boundary separation hypothesis is
\cite[Lemma~4.1]{AJ}. Applying the lemma to the Gromov boundary
therefore gives the form of selective averaging used in the proof.

\section*{Preparation}

Internal calculations and finite reduced-word checks were used during
preparation. The bibliography identifies the versions of the cited works
used in the argument.

\begin{thebibliography}{FOT26}

\bibitem[AD19]{AD}
C. R. Abbott and F. Dahmani,
\emph{Property $P_{\mathrm{naive}}$ for acylindrically hyperbolic groups},
Math. Z. \textbf{291} (2019), 555--568.
\url{https://arxiv.org/abs/1610.04143}.

\bibitem[AJ26]{AJ}
T. Amrutam and Y. Jiang,
\emph{Invariant $C^*$-subalgebras of the reduced group $C^*$-algebra},
arXiv:2503.09548v3, 2026.
\url{https://arxiv.org/html/2503.09548v3}.

\bibitem[FOT26]{FOT}
P. Fima, H. Oppelmayer, and L. Tao,
\emph{Random $C^*$-subalgebras}, arXiv:2610.28295v1, 2026.
\url{https://arxiv.org/html/2610.28295v1}.

\end{thebibliography}
\end{document}
